\section{Introduction}
\label{sec:introduction}

Polynomial graph filters are computationally attractive, but their coefficients do not directly encode localized rational phase geometry. Cayley rational bases provide an alternative spectral coordinate system, yet a fixed basis leaves the location of its pole structure implicit in learned linear coefficients. This work asks whether learnable all-pass roots can provide an interpretable spectral parameterization without sacrificing sparse graph computation.

We study finite Blaschke products composed with a Cayley transform of a symmetric normalized graph Laplacian. The resulting operator is all-pass: it rotates graph spectral phase but cannot attenuate an eigenmode. Spectral selection therefore comes from interference with the identity. The half-sum and half-difference of the all-pass operator form complementary channels whose power responses add to one. Iterating the split emits a sequence of coefficient bands while carrying a complementary remainder.

This construction is inspired by Blaschke phase unwinding, but it is not a graph implementation of classical phase unwinding or a graph analytic signal. An undirected Laplacian does not distinguish the two Fourier orientations in a repeated cycle eigenspace. We therefore frame the method as a learned Blaschke--Cayley tight graph filter bank. The product-sum variant is BDN-inspired because it combines cumulative all-pass factors, while the relaxed GBDN+ model is a separate empirical mixture.

This work makes three contributions. First, we introduce a graph spectral parameterization based on learned Blaschke roots and characterize its phase response and mapped pole geometry. Second, we construct complementary analysis channels and prove that their exact multilevel coefficient map is isometric, with an adjoint synthesis recursion that recovers the analyzed signal. Third, controlled studies distinguish magnitude selection from phase-sensitive recovery and show that mapped-pole distance predicts Chebyshev error more precisely than root radius alone. We additionally report a legacy H100 rerun as preliminary diagnostic evidence. We make no claim about state-of-the-art accuracy, oversquashing, or long-range reasoning before those phenomena are directly evaluated.
